% ============================================================
% Composta con la classe Sapthesis (tesi ufficiali Sapienza).
% ============================================================
\documentclass[binding=0.6cm,italian,oneside,noexaminfo]{sapthesis}

\usepackage[utf8]{inputenc}
\usepackage[T1]{fontenc}
\usepackage[italian]{babel}
\usepackage{graphicx}
\usepackage{xcolor}
\usepackage{booktabs}
\usepackage{longtable}
\usepackage{array}
\usepackage{enumitem}
\usepackage[skip=0.5\baselineskip]{parskip}
\setlist[itemize]{topsep=3pt,itemsep=2pt,parsep=0pt}
\usepackage{listings}
\usepackage{amsmath}
\usepackage{tikz}
\usetikzlibrary{arrows.meta,positioning,fit,backgrounds}
\usepackage{hyperref}

\hypersetup{
    colorlinks=true,
    linkcolor=black,
    citecolor=black,
    urlcolor=black,
    bookmarksopen=true,
}

\definecolor{codebg}{RGB}{245,245,245}
\definecolor{codekw}{RGB}{0,0,150}
\definecolor{codecomment}{RGB}{100,100,100}
\definecolor{codestring}{RGB}{150,0,0}

\lstdefinestyle{codestyle}{
    backgroundcolor=\color{codebg},
    commentstyle=\color{codecomment}\itshape,
    keywordstyle=\color{codekw}\bfseries,
    stringstyle=\color{codestring},
    basicstyle=\ttfamily\small,
    breaklines=true,
    numbers=left,
    numberstyle=\tiny\color{gray},
    frame=single,
    tabsize=2,
    showstringspaces=false,
}
\lstset{style=codestyle}

% ---- Stili per i diagrammi TikZ delle pipeline ----
\tikzset{
  stage/.style={draw, rounded corners=2pt, align=center, font=\footnotesize,
                minimum height=8mm, inner sep=3pt, fill=codebg},
  hl/.style={draw, rounded corners=2pt, align=center, font=\footnotesize,
             minimum height=8mm, inner sep=3pt, fill=blue!8},
  flow/.style={-{Latex[length=2mm]}, thick},
}

% ============================================================
% FRONTESPIZIO (classe Sapthesis)
% ============================================================
\title{Intelligent Packets: inferenza neurale nel data plane eBPF/XDP}
\author{Ettore Cantile}
\IDnumber{2026562}
\course{Engineering in Computer Science}
\courseorganizer{\resizebox{\linewidth}{!}{Dipartimento di ingegneria informatica automatica e gestionale "Antonio Ruberti"}}
\AcademicYear{2025/2026}
\copyyear{2026}
\advisor{Prof. Marco Polverini}
\authoremail{cantile.2026562@studenti.uniroma1.it}
\examdate{}

\begin{document}

\setlength{\parindent}{0pt}

\frontmatter
\maketitle

\tableofcontents

\mainmatter

% ============================================================
\chapter{Fondamenti: eBPF, XDP e il paradigma Intelligent Packets}
\label{ch:fondamenti}
% ============================================================

Questa tesi nasce dall'incontro tra eBPF e il paradigma \textbf{Intelligent Packets} (IPA). Il primo permette di eseguire codice arbitrario, ma sicuro, dentro il kernel Linux; il secondo affida al pacchetto un piccolo modello di machine learning, che ogni nodo esegue per decidere l'inoltro senza consultare un controller centralizzato. Il lavoro unisce queste due possibilità: un modello neurale compatto, quantizzato a interi a 8 bit, viene eseguito interamente in un programma XDP, a velocità di linea. Per comprendere le scelte implementative è utile partire dalle tecnologie su cui si basa il progetto.

\section{eBPF: una macchina virtuale sicura dentro il kernel}

\textbf{eBPF} (extended Berkeley Packet Filter) permette di caricare piccoli programmi nel kernel Linux e di eseguirli in risposta a eventi di sistema o di rete, senza scrivere un modulo kernel né ricompilare o riavviare il sistema. Nato come evoluzione del filtro pacchetti BPF, è diventato negli anni una macchina virtuale general-purpose interna al kernel. Dispone di un proprio set di istruzioni, di un modello di memoria e di funzioni ausiliarie (\emph{helper}) con cui i programmi possono interagire in modo controllato con il resto del sistema.

La possibilità di eseguire codice nel kernel richiede garanzie di sicurezza. In eBPF queste derivano da una sandbox rigorosamente controllata. Prima che un programma venga caricato, un componente del kernel chiamato \textbf{verifier} lo analizza staticamente e ne rifiuta qualunque caratteristica potenzialmente pericolosa. Il verifier costruisce il grafo di controllo del programma e simula tutti i possibili percorsi di esecuzione, verificando una serie di proprietà:
\begin{itemize}
    \item che non esistano cicli il cui numero di iterazioni non sia limitabile staticamente (un programma eBPF deve terminare in un numero finito e noto di passi);
    \item che ogni accesso alla memoria sia entro limiti dimostrabili;
    \item che ogni puntatore usato per un accesso sia stato prima validato;
    \item che il programma non superi un budget massimo di complessità, espresso come numero di istruzioni analizzate lungo tutti i percorsi.
\end{itemize}

Solo se tutte queste verifiche passano, il bytecode eBPF viene compilato \emph{just-in-time} (JIT) in istruzioni macchina native e agganciato al punto di hook richiesto, eseguendo poi alla stessa velocità del codice kernel nativo.

Il verifier garantisce la sicurezza di eBPF, ma impone anche vincoli importanti a chi scrive i programmi. Nel corso del lavoro, diverse scelte implementative sono state guidate non tanto dall'eleganza del codice o dalle prestazioni, quanto dalla necessità di ottenere un programma che superasse la verifica.

Due limiti ricorrono nelle implementazioni descritte in questa tesi. Il primo è il divieto di cicli non limitati, che obbliga a \emph{srotolare} a mano ogni iterazione il cui conteggio non sia noto a compile time. Un prodotto scalare tra un vettore e dei pesi, ad esempio, va espresso come una sequenza di moltiplicazioni e somme, anziché con un \texttt{for}. Il secondo è il limite di complessità: un programma troppo grande, con troppi percorsi possibili, viene rifiutato anche se corretto e deve essere suddiviso in programmi più piccoli. A questi si aggiunge il limite di 512 byte di stack per ogni funzione, che diventerà rilevante nel capitolo sull'evoluzione della pipeline hardcoded.

\section{Punti di aggancio, mappe, XDP e tail call}

eBPF può agganciarsi in molti punti diversi del kernel. Un \emph{kprobe} intercetta l'ingresso o l'uscita da una funzione del kernel o da una syscall, ed è tipicamente usato per tracing e osservabilità. I \emph{tracepoint} sono punti di strumentazione stabili predefiniti nel kernel. Per questo lavoro sono centrali le mappe, l'hook XDP e la primitiva di tail call.

Le \textbf{mappe} (\emph{maps}) sono strutture dati che risiedono nel kernel, ma sono accessibili sia dai programmi eBPF sia dallo user space. Servono a conservare lo stato tra esecuzioni successive e a scambiare dati tra kernel e user space, che possono leggere e scrivere le stesse strutture. Un programma eBPF, infatti, non ha memoria persistente propria oltre alle mappe. Ne esistono diversi tipi: array indicizzati da interi, tabelle hash, array per-CPU, in cui ogni CPU dispone di una copia indipendente senza necessità di lock, e mappe speciali come \texttt{BPF\_PROG\_ARRAY}, che contiene riferimenti ad altri programmi eBPF.

A questi tipi appartengono tutte le mappe del progetto, usate per i pesi del modello, lo stato dei link, le statistiche e la tabella di forwarding.

\textbf{XDP} (eXpress Data Path)~\cite{xdp2018} è il punto di aggancio più vicino alla scheda di rete: il programma eBPF viene eseguito nel driver stesso, prima ancora che il kernel costruisca le proprie strutture dati (l'\texttt{sk\_buff}) per il pacchetto in arrivo. Questa posizione lo rende il punto più veloce per ispezionare, scartare o reindirizzare un pacchetto, ed è il motivo per cui è stato scelto per il data plane del progetto. Un programma XDP riceve un puntatore al pacchetto grezzo e restituisce uno di quattro verdetti principali:
\begin{itemize}
    \item \texttt{XDP\_PASS} lascia proseguire il pacchetto verso lo stack di rete;
    \item \texttt{XDP\_DROP} lo distrugge immediatamente;
    \item \texttt{XDP\_TX} lo rispedisce sull'interfaccia da cui è arrivato;
    \item \texttt{XDP\_REDIRECT} lo devia verso un'altra interfaccia.
\end{itemize}

Quest'ultimo, tramite la helper \texttt{bpf\_redirect()}, è il verdetto su cui si basa tutto il forwarding di questo progetto.

La \textbf{tail call} trasferisce il controllo a runtime da un programma eBPF a un altro, tramite \texttt{BPF\_PROG\_ARRAY} e \texttt{bpf\_tail\_call()}. Se il salto riesce, l'esecuzione non torna al chiamante: il comportamento è più vicino a una \texttt{exec} che a una chiamata di funzione. Nel progetto questo meccanismo serve sia a costruire pipeline modulari, suddividendo la logica in programmi indipendenti, sia a superare il limite di complessità del verifier. Un blocco troppo grande può infatti essere separato in parti che, verificate singolarmente, vengono accettate.

In tutto il codice ricorre poi una precauzione che nelle presentazioni del progetto abbiamo chiamato la ``regola d'oro'' del parsing: prima di accedere a un campo di un header, bisogna verificare esplicitamente che il puntatore non superi \texttt{data\_end}, cioè la fine del buffer del pacchetto. Il verifier non sa a priori quanti byte siano disponibili e, senza questo controllo, rifiuta il programma.

\begin{lstlisting}[language=C,caption={La regola d'oro: controllo dei limiti prima di ogni accesso}]
void *data     = (void *)(long)ctx->data;
void *data_end = (void *)(long)ctx->data_end;
struct ethhdr *eth = data;
if ((void *)(eth + 1) > data_end)
    return XDP_PASS;
if (eth->h_proto == bpf_htons(ETH_P_IP)) {
    struct iphdr *ip = (void *)(eth + 1);
    if ((void *)(ip + 1) > data_end)
        return XDP_PASS;
}
\end{lstlisting}

Mappe, XDP, tail call e parsing sicuro costituiscono la base del sistema. Su di essi si appoggia un protocollo applicativo personalizzato, incapsulato in UDP: l'header ``IPA'' segue quello UDP e usa una porta convenzionale. Il protocollo permette di realizzare il sistema di inferenza neurale interamente nel kernel.

\section{Il paradigma Intelligent Packets}
\label{sec:ipa}

Il progetto IPA propone un modo diverso di pensare al router programmabile. Nel modello tradizionale, un control plane centrale calcola le regole di forwarding e le installa nei nodi, che si limitano a eseguirle. Nel paradigma Intelligent Packets, invece, è il pacchetto a portare l'informazione necessaria per essere instradato. Un modello di machine learning compatto, nel nostro caso un percettrone multistrato (MLP) con pesi quantizzati a 8 bit, viene incorporato nell'header del pacchetto oppure registrato nel nodo. A ogni salto, il router esegue una piccola inferenza locale e sceglie autonomamente l'interfaccia di uscita. Obiettivi diversi, come recuperare da un guasto, rispettare una deadline o gestire la congestione, possono così essere affidati a modelli diversi, senza riconfigurare regole o tunnel.

Il punto di partenza è una dimostrazione preliminare del paradigma. In quella fase l'inferenza nel kernel era limitata a un prodotto scalare a quattro termini tra un piccolo input vector (\texttt{model\_id}, TTL del pacchetto, interfaccia di ingresso e dimensione dell'input) e quattro pesi. L'obiettivo era verificare che l'intero ciclo potesse essere eseguito nel data plane. Lo switch usava due mappe: \texttt{model\_cache}, che associa un \texttt{model\_id} ai pesi del modello, e \texttt{fwd\_table}, che associa il risultato dell'inferenza a un'azione di forwarding, composta dall'interfaccia di uscita e dai MAC da riscrivere. Su questa base erano stati confrontati quattro metodi di popolamento delle tabelle, ciascuno dedicato a un aspetto del problema.

Il primo metodo, basato su quantizzazione post-training (PTQ), addestrava il modello in virgola mobile e lo quantizzava in seguito. Il control plane usava i pesi float originali, mentre il kernel usava quelli interi: questo disallineamento intenzionale produceva errori di forwarding misurabili, utili a quantificare la perdita di accuratezza dovuta alla quantizzazione. Il secondo metodo usava invece l'addestramento consapevole della quantizzazione (QAT), direttamente in interi con fake-quantization. L'aritmetica risultava identica nel control plane e nel kernel, con chiavi sempre corrette e nessun errore.

Il terzo metodo popolava la tabella di forwarding in modo reattivo, in stile OpenFlow. La tabella partiva vuota e ogni fallimento di lookup generava un evento verso il control plane, che installava la regola usando la chiave calcolata dal kernel. In questo modo veniva eliminato il rischio di disallineamento. Il quarto metodo dimostrava il paradigma IPA facendo viaggiare il modello \emph{dentro} il pacchetto. All'inizio, sia la cache dei modelli sia la tabella di forwarding erano vuote. Il primo pacchetto di un dato modello portava i pesi nel payload; il nodo li estraeva, popolava la cache e installava la prima regola. Da quel momento poteva inoltrare autonomamente a velocità kernel, senza coinvolgere di nuovo il control plane.

La fase preliminare ha mostrato che il paradigma funziona end-to-end: il modello può viaggiare nel pacchetto, il nodo può acquisirlo al suo arrivo e usarlo per inoltrare in autonomia. Il modello usato, però, era volutamente minimale. Un singolo prodotto scalare a quattro termini non è una rete neurale multistrato e non permette di affrontare il problema di \emph{come} eseguirla in modo efficiente nel datapath. Questa tesi parte da qui: estende l'inferenza a un MLP multistrato quantizzato e confronta diversi modi di eseguirlo nel kernel, valutandone il compromesso tra prestazioni e flessibilità.

L'inferenza neurale nel kernel è stata esplorata anche da altri lavori. Alcuni impiegano reti convoluzionali leggere in eBPF per classificare il traffico su dispositivi edge~\cite{cnnkernel2025}; altri confrontano reti neurali quantizzate e alberi decisionali eseguiti in kernel e in user space per il packet processing~\cite{nntree2024}. Quest'ultimo studio è il precedente più vicino a questa tesi: la rete viene quantizzata in user space, eseguita nel kernel con aritmetica interamente intera e suddivisa in strati collegati tramite \emph{tail call}, come nella pipeline modular descritta nella Sezione~\ref{sec:p3}. La stessa idea è stata sviluppata anche su hardware programmabile dedicato: N3IC~\cite{n3ic2022} compila reti neurali binarizzate nel data plane di SmartNIC, mentre il survey di Zheng et al.~\cite{inmlsurvey2024} raccoglie le soluzioni basate su switch P4, NetFPGA e SmartNIC. Il contributo di questa tesi non è quindi dimostrare \emph{che} l'inferenza neurale nel datapath sia possibile, ma confrontare \emph{modi diversi} di realizzarla mantenendo invariati modello, hardware e nodo.

Le due fasi sono documentate nel repository. Il ramo \texttt{ipa-poc-preliminar} conserva il codice della dimostrazione iniziale. Il punto di ingresso è \texttt{switch\_core.py}. Sono presenti i quattro metodi \texttt{method1\_ptq}, \texttt{method2\_qat}, \texttt{method3\_openflow} e \texttt{method4\_ipa\_demo}, insieme alle mappe \texttt{model\_cache} e \texttt{fwd\_table} e il prodotto scalare a quattro termini nel kernel. La documentazione del ramo ne riporta metodi, risultati sperimentali e dettagli implementativi.

Il lavoro di questa tesi si trova invece sul ramo principale \texttt{main}, che diverge dal precedente di oltre duecento commit e riorganizza il progetto attorno alle tre pipeline. Nessuno dei quattro metodi originali è rimasto come file, perché è cambiata la domanda di ricerca: non più \emph{come popolare le tabelle} e \emph{quanto costa la quantizzazione}, ma \emph{dove vivono i pesi} e \emph{che cosa può cambiare senza ricompilare}. I due rami documentano quindi fasi successive dello stesso progetto, non varianti alternative; il primo rimane consultabile come premessa del secondo.

\section{Ambiente sperimentale}
\label{sec:lab}

Gli esperimenti sono stati condotti su \emph{Kathara}, uno strumento di emulazione di rete in cui ogni nodo della topologia corrisponde a un container Docker leggero. A differenza di una macchina virtuale completa, il container condivide il kernel dell'host: questo permette di eseguire programmi eBPF/XDP nei nodi emulati. All'avvio, ogni nodo monta il filesystem \texttt{debugfs}, abilita l'inoltro IP e avvia il demone di routing FRR con protocollo OSPF.

La topologia scelta è \emph{Germany50}, importata da SNDlib, un repository di topologie reali usato in letteratura. Comprende una dorsale di cinquanta nodi, a cui si aggiungono due host virtuali come sorgente e destinazione. Il grado massimo di un nodo è sei: il modello deve quindi poter scegliere tra sei interfacce di uscita, a cui si aggiunge la classe di scarto nell'argmax finale. Il traffico di prova usa UDP sulla porta 9999, riservata al protocollo IPA.

Durante la validazione è emerso un limite pratico dell'ambiente: il traffico UDP sulla porta 9999 non viene consegnato end-to-end in tutte le condizioni. Un pacchetto diretto al loopback di un nodo distante più di un salto non arriva, mentre viene consegnato se la destinazione è l'indirizzo IP dell'interfaccia direttamente connessa di un vicino. Per questo la verifica sistematica di correttezza usa \texttt{BPF\_PROG\_TEST\_RUN} anziché traffico reale. La primitiva esegue un programma XDP già caricato su un pacchetto sintetico direttamente nel kernel e restituisce il verdetto e una misura della durata, senza attraversare la scheda di rete.

% ============================================================
\chapter{Il design space: tre pipeline per l'inferenza neurale nel datapath}
\label{ch:designspace}
% ============================================================

Il confronto al centro della tesi riguarda tre modi di eseguire \emph{lo stesso} modello neurale nel datapath eBPF/XDP: un MLP con 65 input, due hidden layer da 4 neuroni ciascuno e 7 output. Le tre pipeline, chiamate P1 hardcoded, P2 template e P3 modular sia nel codice sia nel testo, costituiscono il \emph{design space} del lavoro. Si distinguono per quanto il codice sia specializzato sul singolo modello o, al contrario, generico e parametrizzabile a runtime. La specializzazione aumenta la velocità di picco, ma rende più onerosi gli aggiornamenti; la flessibilità introduce invece più tail call, letture di mappa e istruzioni eBPF per pacchetto, riducendo il throughput massimo. Il capitolo esamina queste tre scelte e ne misura i costi.

\section{Architettura comune alle tre pipeline}
\label{sec:common}

Le tre pipeline condividono il protocollo applicativo e la famiglia di modelli. Il pacchetto IPA viaggia in UDP sulla porta 9999 e contiene un header custom con \texttt{model\_id}, \texttt{scale\_factor} e un descrittore delle feature usate, ripreso nel Capitolo~\ref{ch:hardcoded}. Il programma XDP intercetta il pacchetto prima che raggiunga lo stack di rete Linux ed esegue un percettrone multistrato (MLP) quantizzato post-training a interi a 8 bit. Nel kernel l'\emph{aritmetica è interamente intera}: non compaiono operazioni in virgola mobile, sia perché il verifier non le gestisce, sia per mantenere il calcolo deterministico e riproducibile. L'output è un argmax su tante classi quante sono le interfacce di uscita, più una per lo scarto. Nel nostro scenario le classi sono sette: sei interfacce fisiche e \texttt{DROP}.

L'architettura del modello \emph{non} è fissa. La pipeline template può variare a runtime la larghezza degli strati; la modular può variare anche la profondità; la hardcoded, dopo le modifiche descritte nel Capitolo~\ref{ch:hardcoded}, supporta un numero arbitrario di strati. Per confrontarle, però, occorre un modello \emph{comune}: altrimenti non sarebbe possibile distinguere il costo della pipeline da quello di architetture neurali diverse. Tutte le misure comparative del capitolo usano quindi l'MLP $65 \rightarrow 4 \rightarrow 4 \rightarrow 7$, con sessantacinque feature in ingresso, due hidden layer da quattro neuroni e sette classi in uscita. È il modello addestrato per lo scenario di routing studiato e serve come riferimento di misura, non come vincolo dell'implementazione.

L'input vector contiene 65 componenti, ottenute concatenando quattro tipi di feature legate al routing. Le prime 6 (\texttt{link\_state}) rappresentano lo stato up/down delle sei interfacce di uscita; le successive 6 (\texttt{iface}, codificata one-hot) indicano la porta di ingresso. La tredicesima componente è il TTL del pacchetto, letto come scalare; le ultime 52 (\texttt{node}, anch'essa one-hot) identificano il nodo nella topologia. A runtime solo poche posizioni sono diverse da zero. Tutte e tre le pipeline sfruttano questa sparsità costruendo l'input feature per feature, senza materializzare un array denso di 65 elementi che occuperebbe buona parte dello stack disponibile.

La feature \texttt{link\_state} ha richiesto un intervento specifico. Nella prima versione era sempre a 1, come se tutti i link fossero attivi: il modello non riceveva quindi alcuna informazione sui guasti reali e non poteva reagire. È stato aggiunto il thread in user space \texttt{link\_state\_monitor.py}, che legge periodicamente il file \texttt{carrier} dell'interfaccia \texttt{eth$X$} in sysfs e aggiorna la mappa BPF condivisa \texttt{link\_state[0..5]}.

I programmi eBPF leggono la mappa a ogni pacchetto, senza essere ricaricati. Quando un'interfaccia cade, la feature passa a 0 al polling successivo, eseguito ogni mezzo secondo. L'argmax può così scegliere un'altra uscita: è il meccanismo indicato nella documentazione come \emph{fast-reroute}. All'avvio la mappa viene inizializzata assumendo tutti i link attivi.

Lo schema di elaborazione comune a tutte e tre le pipeline è riassunto nella Figura~\ref{fig:common}: cambia soltanto il blocco centrale, cioè \emph{come} l'inferenza viene eseguita, non il parsing a monte né l'azione a valle.

\begin{figure}[h]
\centering
\begin{tikzpicture}[node distance=6mm]
\node[stage] (in)   {UDP:9999\\header IPA};
\node[stage, right=of in]   (xdp)  {hook XDP\\(parse)};
\node[hl,    right=of xdp]  (inf)  {inferenza\\MLP int8};
\node[stage, right=of inf]  (arg)  {argmax\\7 classi};
\node[stage, right=of arg]  (mac)  {\texttt{mac\_table}\\+ MAC};
\node[stage, right=of mac]  (act)  {\texttt{bpf\_redirect}\\/ \texttt{DROP}};
\draw[flow] (in)--(xdp); \draw[flow] (xdp)--(inf); \draw[flow] (inf)--(arg);
\draw[flow] (arg)--(mac); \draw[flow] (mac)--(act);
\node[draw, dashed, rounded corners, fit=(inf), inner sep=5pt, label=below:{\footnotesize varia tra P1/P2/P3}] {};
\end{tikzpicture}
\caption{Il flusso del pacchetto comune alle tre pipeline. Il solo blocco di inferenza (tratteggiato) cambia implementazione; parsing e azione di forwarding sono identici.}
\label{fig:common}
\end{figure}

\section{Pipeline 1: hardcoded}
\label{sec:p1}

Nella \textbf{pipeline hardcoded} i pesi sono \emph{letterali} inseriti direttamente nel sorgente C, generato e compilato per il singolo modello. Non serve una mappa dei pesi: l'inferenza è una sequenza di operazioni aritmetiche su costanti note. Un dispatcher leggero legge il \texttt{model\_id} dal pacchetto e compie un'unica tail call verso il programma \texttt{model\_<id>} corrispondente. Da quel punto, la costruzione dell'input vector, i due hidden layer, l'output layer e l'argmax sono completamente srotolati, senza cicli né letture di mappa per i pesi. Rimangono soltanto le letture di \texttt{link\_state}, che fornisce una feature di input, e quella di \texttt{mac\_table}, usata dopo l'argmax per determinare l'azione di forwarding.

Conoscere i pesi in fase di compilazione permette di applicare automaticamente la \emph{strength reduction}. Le moltiplicazioni per zero vengono eliminate, quelle per potenze di due diventano shift e altre operazioni vengono semplificate già a compile time. I pesi quantizzati contengono spesso molti zeri e valori piccoli, perciò questa ottimizzazione riduce sensibilmente le istruzioni eseguite per pacchetto. Le altre due pipeline non possono sfruttarla allo stesso modo, perché leggono i pesi da una mappa a runtime.

La gestione delle feature one-hot \texttt{iface} e \texttt{node} è stata più delicata. La soluzione più immediata sarebbe usare un array statico di pesi e indicizzarlo con la posizione attiva, calcolata a runtime, ma questa strada non è accettata dal verifier. Un \texttt{static const} indicizzato da una variabile non risolvibile staticamente genera, al caricamento, una relocation che il loader non sa risolvere e che finisce silenziosamente a zero. Il modello si comporterebbe così come se la feature non contribuisse al risultato, senza che il verifier segnali l'anomalia.

La soluzione adottata usa un solo blocco \texttt{switch} per ciascuna feature one-hot, uno per \texttt{iface} e uno per \texttt{node}. In base alla posizione attiva, il blocco assegna insieme i contributi della feature a \emph{tutti} i neuroni del primo hidden layer. Scrivere uno switch per neurone sarebbe più immediato, ma produrrebbe un'esplosione combinatoria di rami, dell'ordine di $(7 \cdot 52)^{\text{neuroni}}$ percorsi, fino a superare il limite di complessità del verifier. Selezionare insieme i contributi di tutti i neuroni riduce invece il costo a $O(7 + 52)$ rami e permette di superare la verifica.

Il codice~\ref{lst:fc1} mostra il calcolo di un neurone del primo hidden layer. Le componenti \texttt{link\_state} sono lette da una mappa (feature di input); i contributi delle due feature one-hot arrivano dalle variabili \texttt{w\_iface\_j} e \texttt{w\_node\_j} preselezionate dai rispettivi switch; il TTL entra come scalare; tutti i pesi restanti sono letterali cablati nel codice, e la funzione \texttt{RELU\_LL} applica l'attivazione mantenendo l'aritmetica interamente intera a 64 bit per evitare overflow durante l'accumulo.

\begin{lstlisting}[language=C,caption={Un neurone del primo hidden layer, generato per la pipeline hardcoded},label={lst:fc1}]
/* link_state[0..5] letti dalla mappa (feature, non pesi) */
long long h1_0 = RELU_LL(
      (__s64)_ttl * 12LL              /* TTL, scalare, peso letterale */
    + w_iface_0 + w_node_0            /* one-hot: contributi preselezionati */
    + ls0*-7LL + ls1*3LL + ls2*0LL    /* link_state * pesi letterali */
    + ls3*5LL  + ls4*-2LL + ls5*1LL
    + 4LL );                          /* bias, letterale */
\end{lstlisting}
Il termine \texttt{ls2*0LL} contiene un peso nullo, che il compilatore elimina del tutto grazie alla strength reduction, come farà anche con le moltiplicazioni per potenze di due. Nessuna delle altre due pipeline, che leggono i pesi da una mappa a runtime, può beneficiare di questa semplificazione a compile time, perché per il compilatore quei pesi sono valori sconosciuti fino all'esecuzione.

Per aggiornare un modello occorre rigenerare il sorgente C, ricompilarlo e ricaricare l'intero programma XDP. Poiché i pesi sono letterali nel codice, non è possibile modificarli in modo incrementale. La specializzazione offre le prestazioni più alte, come mostrano i risultati sperimentali, ma rende costosi gli aggiornamenti. Il Capitolo~\ref{ch:hardcoded} affronta proprio questo limite.

\section{Pipeline 2: template}
\label{sec:p2}

La \textbf{pipeline template} separa i pesi dal sorgente e li conserva in una mappa BPF (\texttt{arch\_weights}, un \texttt{BPF\_ARRAY}). Il \emph{codice} resta invece fisso e serve un'intera famiglia di architetture. I modelli con due hidden layer e con dimensioni di input e output fissate dal protocollo condividono il programma \texttt{arch\_generic\_2layer}, caricato una sola volta. Le larghezze degli hidden layer, \texttt{n\_h1} e \texttt{n\_h2}, vengono lette a runtime dalla mappa di registro \texttt{arch\_registry}, entro un massimo stabilito in compilazione. Cambiare la larghezza o registrare un altro modello con la stessa profondità richiede quindi solo di scrivere le mappe dei pesi e del registro, senza ricompilare.

Nel primo hidden layer, il prodotto scalare usa un ciclo srotolato anziché sommare termini letterali come in P1. A ogni iterazione legge un peso da \texttt{arch\_weights}, con un offset calcolato a partire da \texttt{n\_h1} e \texttt{n\_h2}. L'intero blocco di pesi occupa una sola entry a valore strutturato: il programma la richiede al kernel una volta, poi accede ai singoli pesi per indirizzamento diretto. La Sezione~\ref{sec:blocco-pesi} misura l'effetto di questa scelta sul costo per pacchetto.

Perché qui il verifier accetta un indice calcolato a runtime, mentre lo rifiuta per l'array statico di P1? \texttt{arch\_weights} è una mappa BPF registrata, cioè un oggetto kernel con dimensioni note e controllate dal verifier, non una relocation da risolvere a compile time. Gli accessi vengono verificati rispetto ai limiti della mappa, senza puntatori irrisolti né esplosioni di rami. Spostare i pesi in una mappa rende quindi P2 più flessibile, ma comporta due costi: un'indirezione per ogni peso e la rinuncia alla \emph{strength reduction}, perché il compilatore non conosce i valori prima dell'esecuzione.

Nel primo layer, il vettore \texttt{link\_state} viene letto una sola volta prima del ciclo sui neuroni e riutilizzato per tutti. Il contributo della componente $i$ al neurone $j$ è \texttt{link\_state[i]} moltiplicato per il peso corrispondente. Per aggiornare un modello o cambiarne la larghezza basta sovrascrivere i pesi e la voce di registro, senza ricompilare né ricaricare il programma.

Questa possibilità è stata verificata registrando nello stesso oggetto BPF due modelli con la stessa profondità ma larghezze diverse: quello reale 65-4-4-7 e uno sintetico 65-6-5-7. Per entrambi, la classe restituita coincideva con quella del riferimento indipendente in Python, confermando il funzionamento della registrazione concorrente di architetture diverse.

\section{Pipeline 3: modular}
\label{sec:p3}

La \textbf{pipeline modular} estende P2 rendendo dinamica anche la \emph{profondità} del modello. L'inferenza attraversa una catena di tail call costruita con due programmi generici. Il primo, \texttt{layer\_first}, viene eseguito al primo hop e legge in modo sparso le 65 feature, come in P1 e P2. Il secondo, \texttt{layer\_hidden}, viene richiamato agli hop successivi e calcola un prodotto denso $n_{in} \rightarrow n_{out}$, con dimensioni lette a runtime dal registro \texttt{layer\_registry}/\texttt{layer\_shapes}.

Le attivazioni intermedie passano da un layer al successivo attraverso una mappa di scratch per-CPU (\texttt{PERCPU\_ARRAY}). Ogni CPU dispone del proprio stato, senza contesa né necessità di locking. Prima di avviare la catena, il dispatcher scrive soltanto i metadati necessari: \texttt{model\_id}, scale factor, indice del layer iniziale, interfaccia di ingresso e TTL. Non costruisce l'intero vettore denso.

L'\emph{ultimo} layer viene riconosciuto dai dati del modello, non dallo slot occupato nella tail-call array. Ogni hop confronta il proprio indice di layer, incrementato di uno, con il numero totale di layer letto dal registro. Se i due valori coincidono, calcola l'argmax e inoltra il pacchetto; altrimenti scrive l'attivazione nello scratch e passa all'hop successivo con una tail call. Modelli di profondità diversa possono così condividere la stessa catena: lo slot 0 contiene sempre \texttt{layer\_first}, quelli successivi \texttt{layer\_hidden}, senza bisogno di riconfigurare i collegamenti.

La verifica ha incluso due modelli registrati nello stesso oggetto compilato: quello reale a 3 layer (65-4-4-7) e uno sintetico a 4 layer (65-5-6-4-7), diverso sia per profondità sia per larghezza. Entrambi hanno prodotto il risultato corretto.

Nello sviluppo di questa pipeline è emerso il limite di complessità del verifier discusso nel primo capitolo. Il primo tentativo, con un unico blocco denso per l'intera inferenza, produceva circa 19983 istruzioni, oltre il limite di 4096 imposto dal kernel di laboratorio. Separare \texttt{layer\_first}, che esegue la lettura sparsa iniziale, da \texttt{layer\_hidden}, che ripete il prodotto denso tramite tail call, è stato quindi necessario per mantenere ogni programma sotto il limite.

Il costo di questa flessibilità dipende dalla profondità del modello: il riferimento a tre layer richiede tre tail call per pacchetto, contro una sola in P1 e P2. Come mostrano le misure, questo è il fattore che incide maggiormente sulle prestazioni di P3.

\section{Dall'argmax all'azione: una mappa comune}

Dopo l'argmax, le tre pipeline traducono la classe scelta nell'azione di rete allo stesso modo. Le classi da 0 a 5 corrispondono alle interfacce, mentre la 6 indica \texttt{DROP}. La classe viene usata come chiave nella mappa \texttt{mac\_table}, che restituisce l'indice dell'interfaccia di uscita e i MAC address sorgente e destinazione da scrivere nell'header Ethernet prima del redirect.

La risoluzione richiede una sola lettura di \texttt{mac\_table} in tutte e tre le pipeline. Una versione precedente di P1 usava uno switch sul valore dell'argmax; sostituirlo con la stessa lettura usata in P2 e P3 ha reso identica l'azione di forwarding, in modo che il confronto riguardi soltanto l'inferenza. Non viene aggiunta un'ulteriore validazione dell'output: la decisione della rete neurale viene accettata, come previsto dal paradigma IPA. La Figura~\ref{fig:pipelines} confronta i tre percorsi, mostrando dove risiedono i pesi, quante tail call servono e che cosa occorre aggiornare per cambiare modello.

\begin{figure}[h]
\centering
\begin{tikzpicture}[node distance=6mm and 9mm]

% ------- P1 hardcoded -------
\node[font=\footnotesize\bfseries] (l1) {P1 hardcoded};
\node[stage, right=of l1,
      text width=18mm] (d1)
    {dispatcher};
\node[hl, right=of d1,
      text width=32mm] (m1)
    {\texttt{model\_<id>}\\
     pesi letterali nel C\\
     inferenza srotolata};
\node[stage, right=of m1,
      text width=22mm] (a1)
    {argmax\\
     + forwarding};
\draw[flow] (d1) -- (m1);
\draw[flow] (m1) -- (a1);

% ------- P2 template -------
\node[font=\footnotesize\bfseries,
      below=10mm of l1] (l2) {P2 template};
\node[stage, right=of l2,
      text width=18mm] (d2)
    {dispatcher};
\node[hl, right=of d2,
      text width=32mm] (m2)
    {\texttt{arch\_generic}\\
     pesi in mappa\\
     \texttt{arch\_weights}};
\node[stage, right=of m2,
      text width=22mm] (a2)
    {argmax\\
     + forwarding};
\draw[flow] (d2) -- (m2);
\draw[flow] (m2) -- (a2);

% ------- P3 modular -------
\node[font=\footnotesize\bfseries,
      below=10mm of l2] (l3) {P3 modular};
\node[stage, right=of l3,
      text width=18mm] (d3)
    {dispatcher};
\node[hl, right=of d3,
      text width=28mm] (f3)
    {\texttt{layer\_first}\\
     primo layer};
\node[hl, right=of f3,
      text width=28mm] (h3)
    {\texttt{layer\_hidden}\\
     $N{-}1$ layer};
\node[stage, right=of h3,
      text width=22mm] (a3)
    {argmax\\
     + forwarding};
\draw[flow] (d3) -- (f3);
\draw[flow] (f3) -- (h3);
\draw[flow] (h3) -- (a3);

\node[below=1.5mm of h3,
      font=\scriptsize,
      align=center]
    {$N$ tail call per pacchetto\\
     stato intermedio in mappa \texttt{PERCPU\_ARRAY}};

\end{tikzpicture}
\caption{Le tre pipeline a confronto. P1 incorpora i pesi come letterali nel codice C (aggiornare un modello richiede rigenerare il sorgente e ricompilare); P2 tiene i pesi in una mappa \texttt{arch\_weights} e li legge a runtime da un unico programma generico; P3 frammenta la rete in una catena di tail call fra \texttt{layer\_first} e \texttt{layer\_hidden}, con pesi in mappe dedicate e attivazioni intermedie in una mappa per-CPU, rendendo dinamiche larghezza e profondità.}
\label{fig:pipelines}
\end{figure}

\section{Le tre pipeline a confronto}
\label{sec:confronto}

Prima di confrontare i risultati, è utile leggere le tre pipeline come una \emph{progressione}. Ciascuna rimuove un vincolo della precedente e introduce un costo aggiuntivo per pacchetto. Non sono quindi tre implementazioni indipendenti, ma tre scelte lungo lo stesso compromesso tra specializzazione e flessibilità.

La distinzione riguarda \textbf{che cosa è fisso e che cosa è variabile senza dover ricompilare}. Nella hardcoded, pesi e architettura sono fissati nel codice: anche modificare un solo peso richiede di rigenerare il sorgente, ricompilarlo e ricaricare il programma. In cambio, il compilatore conosce tutti i pesi e può eliminare l'aritmetica inutile attraverso la strength reduction. Nelle altre due pipeline questa possibilità manca, perché i valori diventano noti solo durante l'esecuzione. La hardcoded sostiene quindi il costo della specializzazione in fase di compilazione e ne beneficia a ogni pacchetto.

La template introduce flessibilità separando i pesi dal codice e conservandoli in una mappa. Il costo di aggiornamento passa così dai secondi della ricompilazione ai millisecondi di una scrittura in mappa. Pesi e larghezza degli hidden layer possono cambiare mentre il datapath è attivo, senza ricaricare il programma. Rimane però fissa la \emph{profondità}: il programma è scritto per due hidden layer, e un modello con tre o quattro strati richiederebbe un programma diverso. Rispetto alla hardcoded, l'aggiornamento è più semplice, ma ogni pacchetto paga l'indirezione dei pesi e l'assenza di \emph{strength reduction}.

La modular rende variabile anche la profondità. La rete è suddivisa in layer generici collegati via tail call e il numero di strati viene letto dal registro del modello. Lo stesso oggetto compilato può quindi servire modelli a due, tre o quattro strati senza modifiche al codice. È la soluzione più flessibile di questo schema, perché l'architettura si aggiorna scrivendo le mappe. È anche la più costosa per pacchetto: ogni strato aggiunge una tail call, con il relativo parsing e il passaggio delle attivazioni attraverso lo scratch per-CPU, aumentando le letture di mappa.

La Tabella~\ref{tab:confronto} riassume le differenze. Da sinistra a destra cresce la flessibilità; le righe mostrano quali costi e vincoli accompagnano questo passaggio.

\begin{table}[h]
\centering
\caption{Le tre pipeline lungo gli assi che le distinguono. Da P1 a P3 la flessibilità cresce, e con essa il costo per pacchetto; specularmente cala il costo di aggiornamento del modello.}
\label{tab:confronto}
\setlength{\tabcolsep}{4pt}
\renewcommand{\arraystretch}{1.15}
\resizebox{\textwidth}{!}{%
\begin{tabular}{@{}llll@{}}
\toprule
 & \textbf{P1 hardcoded} & \textbf{P2 template} & \textbf{P3 modular} \\
\midrule
Dove stanno i pesi & letterali nel codice C & mappa \texttt{arch\_weights} & mappa \texttt{layer\_weights} \\
Variabile senza ricompilare & nulla & larghezza degli strati & larghezza \emph{e} profondità \\
Programmi eBPF compilati & uno per modello & uno per l'intera famiglia & due, generici \\
Tail call per pacchetto & 1 & 1 & $N$ (= profondità) \\
Stato intermedio & nessuno (tutto srotolato) & nessuno & scratch per-CPU \\
Strength reduction & sì (pesi noti) & no & no \\
Aggiornamento modello & ricompila + reload ($\sim$1.1 s) & scrittura mappa ($\sim$0.35 ms) & scrittura mappa ($\sim$0.39 ms) \\
Flessibilità & bassa & media & alta \\
Prestazioni per pacchetto & massime & intermedie & minime \\
\bottomrule
\end{tabular}%
}
\end{table}

La tabella mostra che costo per pacchetto e costo di aggiornamento si muovono in direzioni opposte. La hardcoded è la più conveniente per l'elaborazione dei pacchetti e la più onerosa da aggiornare; la modular presenta il compromesso opposto; la template occupa una posizione intermedia. Nessuna pipeline prevale su entrambi gli aspetti. La scelta dipende quindi dallo scenario d'uso, non soltanto dalla qualità dell'implementazione.

Nel caso studiato, i modelli sono noti in anticipo e cambiano di rado. Conta dunque soprattutto il costo per pacchetto, ed è per questa ragione che la pipeline hardcoded risulta la più adatta, come confermano le misure riportate di seguito.

\section{Risultati sperimentali comparativi}
\label{sec:results}

La misura delle tre pipeline usa \texttt{BPF\_PROG\_TEST\_RUN} per le ragioni discusse nella Sezione~\ref{sec:lab}. Il confronto include una quarta configurazione, la \emph{baseline}: un programma con lo stesso parsing del dispatcher e lo stesso redirect finale, ma senza tail call né inferenza. Questo riferimento permette di distinguere il costo del modello neurale da quello del solo framework XDP.

La suite controlla la correttezza preinstallando nella mappa di forwarding le azioni attese. Dopo l'esecuzione verifica che il verdetto sia coerente con un redirect o un drop e che il contatore della classe scelta sia stato incrementato. La classe restituita dal kernel viene sempre confrontata con quella di un riferimento Python indipendente, che replica la stessa aritmetica intera.

La raccolta dei dati della Tabella~\ref{tab:kernel-results} ha richiesto alcune correzioni alla metodologia iniziale. Usare un solo campione di \texttt{BPF\_PROG\_TEST\_RUN} produceva risultati molto instabili: la stessa configurazione, misurata a pochi minuti di distanza, poteva restituire valori diversi anche di un fattore venti, senza relazione con il numero di istruzioni.

In un ambiente condiviso, interruzioni dello scheduler, altri processi e variazioni della frequenza della CPU introducono un rumore \emph{a senso unico}: possono rallentare un campione, ma non portarlo al di sotto del suo costo reale. Per questo viene riportato il \emph{minimo} di più campioni indipendenti, anziché la media o la mediana, seguendo il principio usato da strumenti come \texttt{hyperfine} e Google Benchmark. Gli errori metodologici nei confronti tra sistemi sono discussi da Heiser~\cite{heiser2018}; per XDP, un intervento a USENIX LISA 2021~\cite{lisa21} mostra anche un caso in cui gli strumenti standard sottostimavano di un fattore sessanta l'utilizzo reale della CPU.

Le misure provengono comunque da una macchina virtuale, senza controllo sulla frequenza della CPU, sui C-state o sull'affinità delle interruzioni. I valori assoluti non sono quindi direttamente confrontabili con quelli ottenuti su hardware dedicato. Le conclusioni si basano sul \textbf{confronto relativo} tra pipeline eseguite sullo stesso nodo e nelle stesse condizioni. Dopo aver modificato la suite per ripetere ogni misura sette volte e riportarne il minimo, i risultati sono diventati stabili tra esecuzioni successive.

Per capire meglio che cosa misurano questi risultati e che cosa lasciano fuori, è utile guardare a come vengono valutati sistemi simili nella letteratura sull'\emph{in-network machine learning}. Il survey di Zheng et al.~\cite{inmlsurvey2024} distingue tre aspetti: le prestazioni del \emph{modello}, come accuratezza e F1; quelle del \emph{sistema}, come latenza e throughput; e le \emph{risorse} richieste sul dispositivo, tra cui memoria, entry di tabella e stadi della pipeline occupati.

La Tabella~\ref{tab:kernel-results} si concentra sugli ultimi due aspetti: riporta la latenza per pacchetto, il throughput, le istruzioni eBPF, il codice jited, le letture di mappa e la memoria occupata dalle mappe. L'accuratezza del modello, invece, non è ciò che vogliamo confrontare: le tre pipeline eseguono la stessa aritmetica intera sugli stessi pesi e devono quindi produrre lo stesso risultato. La suite verifica proprio questo, confrontando la classe scelta dal kernel con quella calcolata da un riferimento Python indipendente. L'accuratezza resta dunque un requisito da preservare, non una differenza da misurare tra le pipeline. Per questo non ha una colonna dedicata in tabella: aggiungerla significherebbe ripetere lo stesso dato.

Per le prestazioni di sistema, la metodologia adottata qui differisce da quella usata in altri lavori. Studi come N3IC~\cite{n3ic2022}, che esegue reti neurali quantizzate nella pipeline di SmartNIC, collegano il sistema sotto test a una seconda macchina che genera traffico reale ad alte prestazioni. Prima delle misure verificano che il collo di bottiglia sia il sistema studiato, non l'apparato di misura. Riportano inoltre distribuzioni o percentili della latenza, anziché un solo valore, e confrontano i risultati con una baseline software ottimizzata a mano, specificando le operazioni incluse.

La valutazione di questa tesi usa invece \texttt{BPF\_PROG\_TEST\_RUN} e non attraversa la rete. Non misura quindi la massima frequenza di inoltro senza perdite, il comportamento sotto saturazione o gli effetti multi-core, come la contesa sulle cache-line discussa nella Sezione~\ref{sec:bottleneck}. Permette però di isolare il costo del programma eBPF da quello del driver, della coda di ricezione e dello scheduler. Questa separazione è utile per il \emph{confronto relativo} tra tre modi di eseguire lo stesso modello sullo stesso nodo, che è l'obiettivo del lavoro. Una valutazione con generatore di traffico rimane la prima estensione dell'attività sperimentale.


\begin{table}[h]
\centering
\caption{Confronto sperimentale tra le quattro configurazioni (kernel reale, \texttt{BPF\_PROG\_TEST\_RUN}, minimo su $N$ trial indipendenti). Sono riportate tre esecuzioni complete e indipendenti della suite: le metriche deterministiche una sola volta quando coincidono, la latenza per tutte e tre. Il confronto fra le tre colonne di latenza è esso stesso un risultato metodologico, discusso nel testo.}
\label{tab:kernel-results}
\resizebox{\textwidth}{!}{%
\begin{tabular}{lrrrr}
\toprule
\textbf{Metrica} & \textbf{baseline} & \textbf{P1 hardcoded} & \textbf{P2 template} & \textbf{P3 modular} \\
\midrule
Istruzioni eBPF (xlated), run A & 129 & 997 & 9\,962 & 8\,209 \\
Istruzioni eBPF (xlated), run B e C & 129 & 997 & 9\,916 & 8\,201 \\
Codice jited (byte), run A & 600 & 4\,751 & 48\,640 & 40\,635 \\
Codice jited (byte), run B e C & 600 & 4\,751 & 50\,306 & 40\,544 \\
Tail call / pacchetto & 0 & 1 & 1 & 3 \\
Letture mappa / pacchetto (reali) & 0 & 4.0 & 8.0 & 26.0 \\
Memoria mappe (byte) & 280 & 308 & 3\,960 & 8\,188 \\
\midrule
Latenza minima (ns/pacchetto), run A (7 trial) & 28.0 & 47.0 & 224.0 & 387.0 \\
Latenza minima (ns/pacchetto), run B (7 trial) & 14.0 & 48.0 & 165.0 & 392.0 \\
Latenza minima (ns/pacchetto), run C (15 trial) & 25.0 & 51.0 & 234.0 & 391.0 \\
Latenza p50 (ns/pacchetto), run C & 33.0 & 52.0 & 244.0 & 419.0 \\
Latenza massima (ns/pacchetto), run C & 53.0 & 67.0 & 323.0 & 529.0 \\
Dispersione relativa $(\max-\min)/\min$, run C & 112\,\% & 31\,\% & 38\,\% & 35\,\% \\
Throughput (Mpps, da latenza minima), run C & 40.000 & 19.608 & 4.274 & 2.558 \\
\bottomrule
\end{tabular}%
}
\end{table}

I risultati seguono l'ordine previsto: passando dalla baseline a P1, P2 e P3 crescono istruzioni, codice jited, tail call, letture di mappa e memoria, mentre le prestazioni diminuiscono. Questo ordinamento si ripete in entrambe le esecuzioni, così come tutte le metriche deterministiche.

Il baseline misura il solo parsing e redirect, senza inferenza, e varia tra $14$ e $28$ nanosecondi per pacchetto. La hardcoded restituisce invece $47$, $48$ e $51$ nanosecondi nelle tre esecuzioni. La tail call, il secondo parsing e l'intera rete neurale aggiungono quindi solo qualche decina di nanosecondi. Questo dato aiuta a rispondere al dubbio sul throughput molto alto della hardcoded: il suo costo rimane nello stesso ordine di grandezza del framework XDP, mentre P2 e P3 se ne allontanano di un fattore tra $9$ e $16$. Il calcolo della rete int8 65-4-4-7 completamente srotolata risulta dunque contenuto rispetto al costo fisso di ricevere e reinoltrare il pacchetto.

Tre riscontri sostengono l'attendibilità dei valori. Il primo è la riproducibilità: nelle tre esecuzioni P1 misura $47$, $48$ e $51$ nanosecondi, mentre P3 misura $387$, $392$ e $391$. La variazione è inferiore al dieci per cento per P1 e all'uno e mezzo per cento per P3. Il secondo è il confronto con gli ordini di grandezza riportati in letteratura: con \texttt{BPF\_PROG\_TEST\_RUN}, un programma XDP minimale costa tipicamente tra dieci e venti nanosecondi e un redirect tra venticinque e cinquanta. Il baseline misurato rientra in questo intervallo.

La colonna delle istruzioni richiede attenzione. Dividere le $997$ istruzioni della hardcoded per i suoi $51$ nanosecondi darebbe circa sei istruzioni e mezzo per ciclo, un valore impossibile su un processore x86, il cui massimo pratico è tra tre e quattro. La colonna \emph{istruzioni eBPF}, però, indica la dimensione \emph{statica} del programma, non le istruzioni eseguite da ogni pacchetto. Lo switch della feature one-hot \texttt{node} contiene cinquantadue rami, ma ne percorre solo \emph{uno}; lo stesso accade per i sei rami della feature di interfaccia. Inoltre, la strength reduction elimina i prodotti per zero e trasforma quelli per potenze di due in scorrimenti. Il percorso eseguito rappresenta quindi solo una parte del codice caricato. Le altre pipeline non beneficiano delle stesse semplificazioni, perché i pesi vengono letti a runtime.

Il \emph{rapporto} rispetto al baseline non è abbastanza stabile per essere riportato come risultato, ed è stato rimosso rispetto a una versione precedente del capitolo. Le latenze delle pipeline si riproducono bene nelle tre esecuzioni, mentre il baseline oscilla tra $28$, $14$ e $25$ nanosecondi. Il rapporto della hardcoded sul baseline diventa così $1.68$, $3.43$ e $2.04$. La dispersione relativa chiarisce il problema: lo scarto tra massimo e minimo del baseline è pari al centododici per cento del minimo, contro il trentuno-trentotto per cento delle pipeline.

Con sole $129$ istruzioni, il baseline ha un costo confrontabile con l'overhead per iterazione di \texttt{BPF\_PROG\_TEST\_RUN}. Per stabilizzarne il minimo servono quindi \emph{più} trial. Passando da sette a quindici, il valore risale a $25$ nanosecondi, vicino ai $28$ della prima esecuzione, e i rapporti delle due esecuzioni concordano entro circa il venti per cento. La normalizzazione può dunque dare un'indicazione, ma non ha precisione sufficiente per questo confronto.

Le latenze sono perciò riportate in valore assoluto, insieme alle caratteristiche della macchina. Il baseline resta un \emph{pavimento qualitativo}, utile a mostrare la differenza tra il framework XDP e l'inferenza, non un divisore. Il costo del salto tra programmi viene misurato separatamente con \texttt{shared/test/bench\_tail\-call\_overhead.py}: lo script confronta due programmi con lo stesso parsing e redirect, diversi soltanto per un hop \texttt{PROG\_ARRAY}, isolando il costo del salto senza ricorrere a rapporti.

P2 e P3 hanno un numero statico di istruzioni molto simile, ma prestazioni nettamente diverse. La differenza dipende quasi interamente dalle tail call, tre contro una, e dalle letture di mappa aggiuntive. Anche questo confronto indica che, nelle condizioni misurate, il costo per pacchetto dipende soprattutto dai salti e dagli accessi a mappa, non dal volume del codice aritmetico.
\newpage
\subsection{Quanto del costo della flessibilità è davvero intrinseco?}
\label{sec:blocco-pesi}

Se gli accessi a mappa dominano il costo di P2 e P3, occorre distinguere due aspetti. Una lettura \emph{per ogni peso} è davvero necessaria per rendere il modello aggiornabile, oppure dipende da un \emph{modo particolare} di implementare questa possibilità? Senza separare le due cose si rischia di attribuire alla flessibilità un costo che deriva invece dalla codifica scelta.

Per verificarlo è stata compilata una build strumentata, in cui ogni \texttt{lookup} incrementa un contatore per-CPU. Questa build è separata da quella usata per misurare le prestazioni. Il conteggio restituiva $147$ letture di mappa per pacchetto in P2 e $160$ in P3, contro le $4$ di P1. Nel codice generato, circa $139$ delle $147$ letture di P2 riguardavano un \emph{singolo byte} di peso, ciascuna con una chiamata a \texttt{bpf\_map\_lookup\_elem} e un controllo di puntatore nullo.

Tutti i pesi erano però già contenuti nella stessa mappa. Era quindi possibile applicare la trasformazione ``da $N$ letture a $1$'', descritta anche per i vettori di feature densi. Dichiarando \texttt{arch\_weights} e \texttt{layer\_weights} come mappe a \emph{valore strutturato}, con un'unica entry per l'intero blocco, basta una sola \texttt{lookup}; gli accessi successivi diventano letture dirette sul puntatore, senza costo aggiuntivo per il verifier. L'indice calcolato a runtime viene mascherato con l'ampiezza del blocco meno uno. Poiché l'ampiezza è una potenza di due, il verifier può dedurre che l'indice resta entro il valore della mappa, eliminando anche i circa $139$ controlli espliciti.

La stessa ottimizzazione si applica allo scratch delle attivazioni intermedie di P3. Prima, le attivazioni venivano rilette con una \texttt{lookup} separata \emph{per ogni coppia} di neurone di uscita e ingresso. L'intero vettore occupa però una sola entry per-CPU e può essere letto una volta per hop.

\begin{table}[h]
\centering
\caption{Effetto del blocco pesi a valore strutturato, a inferenza invariata (stessa classe scelta su tutta la suite di verifica, contro il riferimento Python indipendente).}
\label{tab:blocco-pesi}
\begin{tabular}{lrrrr}
\toprule
 & \multicolumn{2}{c}{\textbf{P2 template}} & \multicolumn{2}{c}{\textbf{P3 modular}} \\
\cmidrule(lr){2-3}\cmidrule(lr){4-5}
\textbf{Metrica} & prima & dopo & prima & dopo \\
\midrule
Letture mappa / pacchetto & 147.0 & 8.0 & 160.0 & 26.0 \\
Istruzioni eBPF (xlated) & 16\,988 & 9\,962 & 15\,767 & 8\,209 \\
Codice jited (byte) & 84\,587 & 48\,640 & 76\,376 & 40\,635 \\
Memoria mappe (byte) & 8\,052 & 3\,960 & 16\,884 & 8\,188 \\
Latenza minima (ns/pacchetto) & 349.7 & 224.0 & 590.0 & 387.0 \\
\bottomrule
\end{tabular}
\end{table}

I risultati sono nella Tabella~\ref{tab:blocco-pesi}. Le letture di mappa per pacchetto scendono da $147$ a $8$ per P2 e da $160$ a $26$ per P3; le istruzioni e il codice jited si riducono sensibilmente; anche la memoria occupata dalle mappe cala, perché un array con valori da un byte paga comunque un arrotondamento per entry che una singola entry compatta non paga.

Le latenze ``prima'' e ``dopo'' della Tabella~\ref{tab:blocco-pesi} sono state raccolte nella \emph{stessa sessione, sulla stessa macchina}, modificando soltanto il contenitore dei pesi. I valori assoluti sono quindi direttamente confrontabili, senza normalizzazione: P2 passa da $349.7$ a $224.0$ nanosecondi per pacchetto e P3 da $590.0$ a $387.0$, con un guadagno di circa un terzo per entrambe. Tra tabelle raccolte su macchine diverse, invece, il confronto deve limitarsi alle metriche deterministiche, come istruzioni, letture di mappa e memoria, per le ragioni già discusse sull'instabilità del baseline.

Questo risultato cambia l'interpretazione delle misure precedenti. Circa un terzo del costo di P2 e P3 \emph{non} dipendeva dalla flessibilità, ma dalla codifica del contenitore dei pesi. È stato possibile eliminarlo senza cambiare l'architettura a dispatcher e tail call, la genericità del programma o la possibilità di aggiornare il modello scrivendo una mappa.

L'ordine delle pipeline rimane lo stesso, con P1 più veloce di P2 e P2 più veloce di P3, ma il divario si riduce. Il costo va quindi separato in una componente \emph{strutturale}, dovuta alle tail call e all'indirezione dei pesi letti a runtime, che non si elimina senza rinunciare alla flessibilità, e una componente \emph{implementativa}, che dipende da come vengono organizzati gli accessi. Misurare e rimuovere quest'ultima prima di attribuire un costo al design space è uno dei risultati metodologici del confronto.

Il confronto comprende anche il costo di registrare un nuovo modello a runtime, misurato con chiamate reali a \texttt{BPF\_PROG\_LOAD} e \texttt{bpf\_map\_update\_elem}. È l'altro lato del compromesso: quanto costa aggiornare il modello, anziché eseguirlo. Nella hardcoded occorre rigenerare e ricompilare il sorgente C da zero, con un costo medio di circa 1.1 secondi. Template e modular scrivono invece le nuove mappe dei pesi e richiedono circa 0.35 e 0.39 millisecondi, rispettivamente: oltre tremila volte meno.

Anche per questi aggiornamenti viene riportato il minimo, non la media, per il carattere unidirezionale del rumore già discusso. La prima registrazione paga l'accesso iniziale alle pagine della mappa appena creata e può dominare la media di poche misure, senza rappresentare il costo dell'operazione a regime. Il blocco pesi a valore strutturato della Sezione~\ref{sec:blocco-pesi} riduce anche questo costo: ora il modello viene registrato con una sola \texttt{bpf\_map\_update\_elem} sull'intero blocco, anziché una chiamata per peso. Template e modular scambiano quindi un costo per pacchetto più alto con aggiornamenti molto più rapidi.

La suite non si limita al modello di riferimento 65-4-4-7, ma verifica correttezza e prestazioni anche su forme diverse. Per la hardcoded vengono compilati e controllati modelli con un solo hidden layer e con tre hidden layer. Per template e modular, un modello sintetico con forma diversa viene registrato insieme a quello reale nello stesso oggetto compilato. Tutte le varianti superano il confronto di correttezza con il riferimento Python.

\section{Colli di bottiglia trasversali del datapath}
\label{sec:bottleneck}

Le differenze viste finora dipendono dal lavoro svolto da ciascuna pipeline per pacchetto: tail call, letture di mappa e istruzioni. Alcuni costi sono però comuni a tutte e tre e non dipendono da come viene eseguita l'inferenza. Analizzarli permette di individuare ottimizzazioni applicabili all'intero datapath. I quattro punti seguenti sono stati riscontrati nel codice C generato per P1, P2 e P3.

Il primo è il \textbf{doppio parsing} del pacchetto. Il dispatcher, prima di scegliere il programma da invocare, esegue il parsing completo degli header, Ethernet, IP, UDP e infine l'header IPA, solo per poter leggere un singolo byte, il \texttt{model\_id}, e decidere verso quale \texttt{model\_<id>} fare tail call:
\begin{lstlisting}[language=C,caption={Il dispatcher parsa tutto solo per leggere \texttt{model\_id}}]
/* ... controlli su eth, ip, ip_proto, ihl, udp->dest == 9999 ... */
struct ipa_hdr *ipa = (struct ipa_hdr *)(udp + 1);
if ((void *)(ipa + 1) > data_end) return XDP_PASS;
__u32 mid = (__u32)ipa->model_id;
model_progs.call(ctx, mid);   /* tail call: i puntatori qui muoiono */
\end{lstlisting}
Dopo la tail call, il programma \texttt{model\_<id>} deve ripetere lo stesso parsing. I puntatori \texttt{data}, \texttt{data\_end} e quelli derivati non sopravvivono infatti al salto. Il programma riceve il contesto \texttt{ctx}, da cui deve ricavare nuovamente \texttt{data}/\texttt{data\_end} e ripetere i controlli di limite; altrimenti il verifier rifiuta l'accesso al pacchetto. Ogni pacchetto attraversa quindi due volte il parsing degli header, contribuendo ai $\sim$14 nanosecondi di differenza tra hardcoded e baseline.

Nel caso a modello singolo si potrebbe agganciare direttamente il programma del modello, eliminando dispatcher, tail call e secondo parsing. Si perderebbe la possibilità di dispatch multi-modello, ma sarebbe un compromesso ragionevole per la hardcoded quando il modello è noto in anticipo.

Il secondo è la scelta del tipo di mappa per \texttt{mac\_table}, che traduce la classe scelta dall'argmax nell'azione di forwarding. Nella versione iniziale delle tre pipeline, questa mappa era dichiarata come \texttt{BPF\_HASH}, con chiave un intero che rappresenta la classe (da 0 a 6):
\begin{lstlisting}[language=C,caption={\texttt{mac\_table} come hash, con chiavi intere dense}]
BPF_HASH(mac_table, __u32, struct fwd_action, 8);
/* ... */
__u32 _cls = (__u32)best_cls;   /* 0..6 */
struct fwd_action *_action = mac_table.lookup(&_cls);
\end{lstlisting}
Una tabella hash calcola l'hash della chiave e gestisce le eventuali collisioni percorrendo un bucket. Questo costo è giustificato per chiavi sparse o non note, mentre qui le chiavi sono gli interi consecutivi da 0 a 6. Una \texttt{BPF\_ARRAY}, indicizzata dalla classe, può svolgere lo stesso compito con un accesso a offset costante, senza hash né collisioni. La modifica richiede poco lavoro e si applica allo stesso modo alle tre pipeline.

Il terzo punto riguarda i contatori di statistica e può incidere soprattutto su hardware multi-core reale. Le mappe \texttt{pkt\_stats} e \texttt{cls\_stats} sono \texttt{BPF\_ARRAY} condivise tra tutte le CPU, e vengono aggiornate ad ogni pacchetto con un'operazione atomica:
\newpage
\begin{lstlisting}[language=C,caption={Contatori condivisi aggiornati con atomica per pacchetto}]
int _hi = 0; __u64 *_hv = pkt_stats.lookup(&_hi);
if (_hv) __sync_fetch_and_add(_hv, 1);   /* atomica su cache-line condivisa */
__u64 *_cv = cls_stats.lookup(&_cls);
if (_cv) __sync_fetch_and_add(_cv, 1);
\end{lstlisting}
Su hardware reale, RSS (Receive Side Scaling) distribuisce il traffico XDP tra più CPU. Ogni CPU aggiorna però i contatori sulla \emph{stessa} riga di cache, che passa continuamente da un core all'altro: è il fenomeno del \emph{cache-line bouncing}. Il costo cresce con il numero di CPU e non emerge nelle misure su singola CPU con \texttt{BPF\_PROG\_TEST\_RUN}.

Usando \texttt{BPF\_PERCPU\_ARRAY}, ogni CPU può aggiornare una copia privata senza atomiche né contesa; lo user space somma le copie quando legge le statistiche. La modular usa già mappe per-CPU per le attivazioni intermedie. Estendere la stessa soluzione ai contatori di tutte le pipeline sarebbe quindi coerente con il codice esistente e a basso rischio.

Per il redirect finale, tutte le pipeline usano \texttt{bpf\_redirect(ifindex, 0)}, passando direttamente l'indice dell'interfaccia di uscita.

Si può usare invece \texttt{bpf\_redirect\_map()}, che legge una voce di una mappa \texttt{DEVMAP}. Oltre a essere leggermente più veloce sul singolo pacchetto, permette l'\emph{invio in blocco}: i pacchetti vengono accumulati e trasmessi insieme alla fine del ciclo di polling della scheda, anziché uno alla volta. La differenza conta soprattutto ad alti volumi di traffico ed è il motivo per cui la documentazione XDP raccomanda il percorso \texttt{DEVMAP} nei data plane di produzione.

La conversione di \texttt{mac\_table} da \texttt{BPF\_HASH} a \texttt{BPF\_ARRAY} è stata applicata a tutte e tre le pipeline. È stato necessario adattare anche il riconoscimento delle classi non configurate: poiché una \texttt{BPF\_ARRAY} non restituisce un puntatore nullo, questa condizione viene ora segnalata da un \texttt{ifindex} pari a zero, che non identifica un'interfaccia di uscita valida. I dati della Tabella~\ref{tab:kernel-results} sono stati raccolti dopo la modifica e ne includono gli effetti.

Gli altri tre interventi, eliminazione del doppio parsing, contatori per-CPU e redirect tramite \texttt{DEVMAP}, non sono stati applicati. Richiederebbero una nuova campagna di misure sull'hardware di validazione; alcuni effetti, soprattutto la contesa dei contatori, si osservano solo su hardware multi-core reale e non nel laboratorio virtualizzato. Restano comunque ottimizzazioni indipendenti dalla scelta tra hardcoded, template e modular: riguardano il framework comune e sono possibili sviluppi di una futura versione del datapath.

\section{Il costo della flessibilità e la scelta della pipeline vincente}

La scelta tra le pipeline dipende da quanto si vuole specializzare il codice sul singolo modello:
\begin{itemize}
    \item La pipeline hardcoded massimizza le prestazioni a scapito della flessibilità;
    \item la template trova un compromesso pratico, un solo programma per un'intera famiglia di architetture aggiornabile scrivendo una mappa;
    \item la modular massimizza la flessibilità, rendendo dinamiche larghezza e profondità, al prezzo del numero di tail call più alto.
\end{itemize}

Le misure indicano quindi la hardcoded come la soluzione più adatta allo scenario studiato, in cui i modelli sono noti in anticipo e cambiano di rado. Per questo il resto della tesi si concentra su di essa. I suoi due limiti principali, l'input vector legato a un solo scenario e il costo elevato di aggiornamento, dipendono da scelte implementative specifiche, non dalla specializzazione in sé. Il capitolo seguente mostra come affrontarli senza rinunciare al vantaggio prestazionale.

% ============================================================
\chapter{L'evoluzione della pipeline hardcoded}
\label{ch:hardcoded}
% ============================================================

La pipeline hardcoded aggiunge appena 14 nanosecondi al costo del solo framework XDP e offre le prestazioni migliori per lo scenario considerato. Questo risultato ha motivato il lavoro sui due limiti durante la revisione: l'input vector era legato a un solo scenario di rete e ogni aggiornamento richiedeva una ricompilazione costosa. Il capitolo descrive come sono stati affrontati entrambi. Presenta inoltre un terzo filone, emerso durante lo sviluppo: capire se, a parità di budget di pesi, convenga costruire un modello più largo o più profondo.

Per tutto il lavoro è rimasta ferma una distinzione: nello scenario con modelli noti a priori, la dimensione di ciascun tipo di feature dipende dalla \emph{topologia} ed è condivisa dai nodi della stessa rete. Non è un parametro scelto liberamente dal singolo modello. Ogni modello decide \emph{quali} tipi di feature usare, non la loro dimensione. Questa separazione permette di rendere generico l'input senza perdere i vantaggi della specializzazione.

\section{Un input vector generico, senza perdere prestazioni}

La versione originale prevedeva la forma 65-x-x-7, con sessantacinque feature in ingresso nello stesso ordine: \texttt{link\_state}, \texttt{iface}, \texttt{ttl} e \texttt{node}. Questa configurazione era pensata per il routing basato su FRR. Per usare feature diverse, o cambiare il numero di componenti di una feature, occorreva riscrivere a mano il generatore di codice.

La soluzione è stata introdurre un \emph{catalogo} dei tipi di feature che il generatore sa costruire, e un \textbf{descrittore per-modello} che elenca, in ordine, quali tipi quel modello usa. Il catalogo distingue tre categorie (\emph{kind}) di feature:
\begin{itemize}
    \item \emph{scalar}, un singolo valore letto dal pacchetto in transito, come il TTL;
    \item \emph{dense\_vector}, un piccolo vettore letto in un colpo solo da una mappa BPF condivisa, come \texttt{link\_state} o una nuova feature \texttt{queue\_occupancy} aggiunta in questa fase;
    \item \emph{onehot}, esattamente un indice attivo su un range di dimensione nota, come \texttt{iface} o \texttt{node}.
\end{itemize}

Il descrittore del modello è una lista ordinata di tipi. La \emph{dimensione} di ciascuno viene invece letta da \texttt{topology\_config}, un file condiviso da tutti i nodi della rete: da qui si ricavano, ad esempio, il numero di interfacce di \texttt{link\_state} e il numero di nodi di \texttt{node}. Le eventuali chiavi topologiche presenti nel file del modello vengono ignorate, perché la topologia appartiene alla rete, non al singolo modello.

\begin{lstlisting}[language=Python,caption={Descrittore per-modello: solo i TIPI, non le dimensioni}]
features = ["link_state", "queue_occupancy", "ttl"]
hidden_dims = [4, 4]
n_out = 7
# N_IN = somma delle dimensioni dei tipi elencati,
#        lette dal topology_config della rete
\end{lstlisting}

Per ogni tipo elencato nel descrittore, il generatore emette il codice corrispondente al kind. Una feature scalare viene letta dal pacchetto; una dense richiede una lettura di mappa seguita dagli accessi diretti al valore restituito; una one-hot usa il blocco \texttt{switch} descritto nel capitolo precedente. L'input viene così costruito feature per feature, mantenendo l'inferenza completamente srotolata.

Se il modello non specifica un descrittore, il generatore usa la configurazione storica \texttt{[link\_state, iface, ttl, node]}, preservando il comportamento precedente e il programma byte-identico nel caso storico. Il confronto riga per riga tra vecchio e nuovo generatore, a parità di pesi, ha rilevato soltanto differenze cosmetiche nell'ordine dei termini della somma, senza differenze sostanziali.

La maggiore flessibilità introduce un rischio: un modello può essere addestrato per una dimensione di input diversa da quella prodotta dalla topologia del nodo. Inizialmente il problema passava inosservato. Il programma veniva generato e caricato senza errori, ma applicava i pesi a un layout di input incompatibile e produceva un'inferenza sbagliata.

È stato quindi aggiunto un controllo prima della generazione del codice. La dimensione attesa viene letta dal primo strato lineare del checkpoint e confrontata con quella ricavata dalla topologia per il descrittore scelto. Se non coincidono, viene sollevata un'eccezione. Il caricamento si interrompe così con un errore esplicito, anziché lasciare in esecuzione un modello che sceglie next-hop errati senza segnalarlo.

La generalizzazione dell'input ha permesso anche di ottimizzare il percorso di elaborazione dei pacchetti in tutte e tre le pipeline. Nella prima versione, ogni componente di una feature dense richiedeva una chiamata a \texttt{bpf\_map\_lookup\_elem()}: sei per \texttt{link\_state} e quattro per \texttt{queue\_occupancy}, fino a dieci helper call per pacchetto soltanto per queste feature. I valori di ogni feature possono però essere raccolti in una struttura che occupa un'unica entry della mappa. Basta allora una sola \texttt{lookup} su chiave fissa, seguita da letture dirette sul puntatore restituito, senza il costo aggiuntivo delle helper call per il verifier.

\newpage
\begin{lstlisting}[language=C,caption={Da N lookup a 1: mappa struct-valued}]
struct ls_vec { __u32 v[6]; };
BPF_ARRAY(link_state, struct ls_vec, 1);

int z = 0;
struct ls_vec *p = link_state.lookup(&z);
if (p) { ls0 = p->v[0]; ls1 = p->v[1]; /* ... */ }
\end{lstlisting}

La modifica è stata applicata a tutte e tre le pipeline senza cambiare l'architettura a dispatcher e tail call. Nella hardcoded riduce di cinque le letture di mappa per pacchetto e lascia invariata l'inferenza, verificata su tutta la suite confrontando la classe scelta con il riferimento Python. Nei confronti tra versioni, il numero di letture risulta quindi \emph{diminuito} nonostante la maggiore flessibilità dell'input.

\section{AOT-literal: separare la compilazione dal deploy}
\label{sec:aot}

Il costo di ricompilazione richiede un approccio diverso. Nella hardcoded i pesi sono letterali nel codice, quindi modificarli impone una nuova compilazione. L'obiettivo non è eliminarla, ma spostarla \emph{fuori dal percorso critico}.

Nella versione originale, la hardcoded usa BCC (BPF Compiler Collection), una libreria che incorpora LLVM/clang e compila il sorgente C sul nodo datapath a ogni caricamento. La voce \texttt{[M1 update timing]} della suite misura l'intero ciclo di rigenerazione, compilazione e ricaricamento. Il costo è di circa 1.3 secondi per modello, con valori tra 1.26 e 1.66 secondi a seconda della macchina, quasi interamente dovuti alla chiamata a \texttt{clang}. È accettabile per aggiornamenti rari, ma diventa oneroso quando lo stesso modello va distribuito a molti nodi o sostituito sul datapath in esecuzione.

Serve ancora una compilazione per modello, ma possono cambiare \emph{dove} e \emph{quando} viene eseguita. Il percorso alternativo, chiamato \textbf{AOT-literal} (Ahead-Of-Time), genera lo stesso tipo di sorgente C con pesi letterali, usando il dialetto \texttt{libbpf} anziché BCC. La compilazione con \texttt{clang} avviene \emph{offline}, su una build box separata, e produce un file oggetto \texttt{.o} pronto per il caricamento. Al nodo datapath rimane soltanto il compito di aprirlo e caricarlo, con un costo di ordini di grandezza inferiore alla compilazione da zero. La Figura~\ref{fig:aot} mostra questa separazione.

\begin{figure}[h]
\centering
\begin{tikzpicture}
\tikzset{
  bx/.style={draw, rounded corners=2pt, align=center, font=\footnotesize,
             minimum height=9mm, minimum width=30mm, inner sep=3pt, fill=codebg},
  cx/.style={bx, fill=blue!8},
}
\node[cx] (gen)   at (0,0)      {genera C\\(pesi letterali)};
\node[cx] (clang) at (4.2,0)    {\texttt{clang -O2}\\$\rightarrow$ \texttt{.o}};
\node[bx] (open)  at (0,-3.3)   {\texttt{open} + \texttt{load}\\(verify + JIT)};
\node[bx] (att)   at (4.2,-3.3) {attach XDP};
\draw[flow] (gen)--(clang);
\draw[flow] (open)--(att);
\draw[flow] (clang.south) -- ++(0,-1.05) -| (open.north);
\node[font=\scriptsize, anchor=south] at (2.1,-1.45) {\texttt{.o} prebuilt};
\begin{scope}[on background layer]
  \node[draw, dashed, rounded corners, fit=(gen)(clang), inner sep=7pt,
        label=above:{\footnotesize build box $\cdot$ offline $\cdot$ $\sim$1$\times$ per modello}] {};
  \node[draw, dashed, rounded corners, fit=(open)(att), inner sep=7pt,
        label=below:{\footnotesize nodo datapath $\cdot$ $\sim$6\,ms $\cdot$ niente \texttt{clang}}] {};
\end{scope}
\end{tikzpicture}
\caption{AOT-literal: la compilazione \texttt{clang} (costosa) resta una per modello ma è spostata offline sulla build box; sul nodo datapath resta solo il caricamento del file oggetto già compilato, che costa pochi millisecondi contro il $\sim$1.3 secondi della compilazione BCC a runtime.}
\label{fig:aot}
\end{figure}

Le misure mostrano il vantaggio di questa separazione. La compilazione offline con \texttt{clang} richiede tra 90 e 555 millisecondi, a seconda dell'esecuzione, sostenuti una sola volta sulla build box. Sul nodo datapath, apertura del file e caricamento nel kernel, inclusi verifica e JIT, richiedono invece tra 4.7 e 6.2 millisecondi, contro circa 1.3 secondi del percorso BCC per lo stesso modello. Il guadagno riguarda quindi il deploy, non la latenza per pacchetto: il costo si riduce di oltre due ordini di grandezza e \texttt{clang} non deve più essere presente sul nodo.

Le prestazioni per pacchetto non sono invece identiche a quelle del percorso BCC. L'architettura resta la stessa, con dispatcher, tail call e doppio parsing, e i pesi letterali mantengono il beneficio della strength reduction. Cambiano però la versione di \texttt{clang} e il modo di accedere alle mappe: da un lato il riscrittore di BCC, dall'altro \texttt{bpf\_map\_lookup\_elem} di \texttt{libbpf}. Il codice prodotto differisce quindi di qualche punto percentuale.

Nella stessa sessione e con la metodologia della tabella comparativa, BCC produce 997 istruzioni totali, 29 del dispatcher e 968 del modello, con una latenza di 48 nanosecondi per pacchetto. AOT produce 1\,026 istruzioni, 28 più 998, e misura 57 nanosecondi. AOT rimane dunque nella stessa \emph{classe di prestazioni} della hardcoded, ben distante dai 165 e 392 nanosecondi di P2 e P3, ma non genera un programma identico bit per bit. Il miglioramento principale resta il costo di aggiornamento sul nodo, ridotto di oltre due ordini di grandezza.

La prima versione di AOT-literal supportava soltanto il descrittore storico, perché il catalogo di feature era stato implementato nel solo dialetto BCC. I generatori per i tre kind, scalar, dense e one-hot, sono stati quindi portati anche in \texttt{libbpf}. Ogni descrittore accettato dal percorso BCC può ora essere compilato anche in AOT. La differenza riguarda la sintassi, non il funzionamento: alle funzioni come \texttt{m.lookup(\&k)} di BCC corrispondono le dichiarazioni \texttt{struct \{...\} SEC(".maps")} e le chiamate \texttt{bpf\_map\_lookup\_elem(\&m, \&k)} di \texttt{libbpf}.

AOT-literal è diventato quindi l'\emph{unico} percorso di distribuzione e aggancio della hardcoded, sostituendo BCC come backend di deploy. Per farlo è stato necessario risolvere un problema dei nodi Kathara: non hanno né \texttt{clang} né \texttt{libbpf} installati. BCC incorpora LLVM/clang come libreria e non richiede questi strumenti esterni; inizialmente questa differenza aveva motivato il mantenimento di entrambi i percorsi.

AOT non ha mai richiesto di eseguire \texttt{clang} sul nodo: la compilazione avviene sulla build box e viene trasferito solo il file oggetto. Il problema era nel piccolo loader C incaricato di aprirlo e caricarlo nel kernel. Essendo collegato dinamicamente a \texttt{libbpf}, non poteva avviarsi sui nodi privi della libreria condivisa.

Il loader è stato quindi collegato staticamente a \texttt{libbpf} e alle sue dipendenze, \texttt{libelf}, \texttt{zlib} e, sulle distribuzioni recenti, \texttt{libzstd} e \texttt{liblzma}. Rimane dinamico soltanto il collegamento alla libreria C standard, normalmente presente sui sistemi Linux. Si ottiene così un unico eseguibile, compilato una volta sulla build box e utilizzabile su tutti i nodi del laboratorio, privi di \texttt{clang}, \texttt{libbpf} e delle librerie di sviluppo eBPF. Il deploy della hardcoded usa ora soltanto AOT-literal; il percorso BCC non viene più invocato per il caricamento live.

BCC rimane nell'infrastruttura di test, dove serve a compilare e verificare offline la correttezza dell'inferenza. La suite viene eseguita su una macchina di sviluppo con BCC disponibile, non sul nodo datapath in produzione. Il suo utilizzo è quindi separato dal meccanismo di distribuzione del modello.

\section{Profondità variabile e studio larghezza contro profondità}
\label{sec:width-depth}

Sia il generatore BCC sia quello AOT erano stati scritti per modelli con esattamente due hidden layer. Le variabili \texttt{n\_h1} e \texttt{n\_h2}, la suddivisione dei pesi e l'emissione del codice riflettevano questa scelta. I due generatori, insieme al riferimento Python della suite, sono stati generalizzati per accettare una sequenza di hidden layer di lunghezza qualsiasi: due come prima, ma anche uno, quattro o nessuno, nel caso di un modello lineare.

Il layout dei pesi segue ora ogni coppia di strati consecutivi: un blocco di $n_{prev} \times n_{cur}$ pesi e, a seguire, $n_{cur}$ bias. Con due hidden layer si ottiene lo schema precedente. La compatibilità è stata verificata confrontando byte per byte il codice generato, risultato identico a meno delle righe vuote, e confrontando su pesi e input casuali il nuovo riferimento Python con la precedente implementazione esplicita a due strati.

La generalizzazione ha permesso di esplorare una domanda ulteriore: a parità di numero di pesi, e quindi approssimativamente di capacità del modello, conviene usare pochi strati larghi o molti strati stretti? La risposta ha conseguenze pratiche per la hardcoded, perché l'intera inferenza viene srotolata in un'unica funzione C, soggetta al \textbf{limite di 512 byte di stack} imposto dal verifier.

Uno script di benchmark dedicato confronta tre budget crescenti, circa 300, 1200 e 4700 pesi. Per ciascuno considera una forma ``larga'', con un solo hidden layer, e due forme ``profonde'', con quattro e otto hidden layer. I budget non sono perfettamente uguali tra le forme: lo scarto effettivo viene sempre riportato.

Il descrittore storico contiene la feature one-hot \texttt{node}, con 52 valori, che incide sul costo. Per evitare che fosse questa scelta a determinare il risultato, l'esperimento è stato ripetuto con quattro descrittori: quello storico con due feature one-hot, uno senza, uno con una sola one-hot piccola e uno con una sola one-hot grande. Le latenze sono calcolate come minimo su più trial indipendenti, secondo la metodologia della Sezione~\ref{sec:results}, per contenere l'effetto del rumore.

Le configurazioni più grandi hanno posto un problema distinto. Quando superano il limite di stack, il backend LLVM di BCC termina il processo con un errore fatale, non con un'eccezione intercettabile. Prevedere quali forme falliscano si è rivelato inaffidabile: il compilatore riesce a volte a riutilizzare lo stack tra variabili con vita non sovrapposta, ma non in modo prevedibile. Ogni configurazione viene quindi eseguita in un sottoprocesso isolato. Se questo termina con un errore fatale, la configurazione è registrata come fallita e il benchmark prosegue con le altre.

\begin{table}[h]
\centering
\caption{Larghezza contro profondità a parità di budget di pesi (minimo su 15 trial, descrittore storico).}
\label{tab:width-depth}
\resizebox{\textwidth}{!}{%
\begin{tabular}{lrrr}
\toprule
\textbf{Configurazione} & \textbf{Pesi} & \textbf{Latenza minima (ns/pkt)} & \textbf{Note} \\
\midrule
Baseline / largo, 1$\times$4 & $\sim$300--320 & 44--56 & riferimento \\
Profondo, 8$\times$3 & $\sim$310 & 57--76 & overhead contenuto \\
\midrule
Largo, 1$\times$16 & $\sim$1175 & 103--111 & sempre il più veloce \\
Profondo, 4$\times$11 & $\sim$1206 & 203--236 & circa 2$\times$ più lento \\
Profondo, 8$\times$9 & $\sim$1294 & 269--301 & 2.5--3$\times$ più lento \\
\midrule
Qualunque forma & $\sim$4700--4780 & --- & \textbf{sempre fallito} \\
\bottomrule
\end{tabular}%
}
\end{table}

La Tabella~\ref{tab:width-depth} riassume un risultato coerente nei quattro descrittori: a parità di budget di pesi, allargare il modello è sistematicamente più veloce che aggiungere profondità. La differenza cresce con il budget e raggiunge un fattore due o tre al livello intermedio. Ogni hidden layer aggiuntivo introduce un costo fisso di circa 15--40 nanosecondi, dovuto alla transizione tra strati e all'attivazione ReLU, indipendentemente dalla composizione delle feature.

Al budget più alto, \emph{ogni} configurazione supera il limite di stack e fallisce in compilazione, sia larga sia profonda e con qualunque descrittore. L'attuale architettura hardcoded, con inferenza srotolata in una sola funzione e pesi letterali, incontra quindi un limite di capacità che non dipende dalla distribuzione dei pesi tra gli strati. Per superarlo occorre spostare gli array più grandi in una mappa \texttt{PERCPU\_ARRAY}, come suggerisce anche il messaggio del compilatore, anziché continuare a ridistribuire lo stesso numero di pesi.

\section{Nota conclusiva sull'evoluzione}

Gli interventi descritti nel capitolo ampliano le possibilità della hardcoded mantenendone il vantaggio prestazionale nello scenario studiato. L'input vector generico consente di usare il modello in contesti diversi e conserva il codice byte-identico nel caso storico. AOT-literal riduce il costo di aggiornamento di quasi tre ordini di grandezza, mantenendo prestazioni per pacchetto comparabili a quelle del percorso precedente. La profondità variabile permette di confrontare architetture diverse e di individuare sia una scelta pratica, preferire la larghezza alla profondità, sia il limite oltre il quale il design attuale non può crescere.

La pipeline più specializzata, inizialmente anche la più rigida, ha quindi offerto margini di evoluzione concreti.

% ============================================================
\begin{thebibliography}{9}

\bibitem{xdp2018}
T.~Høiland-Jørgensen, J.~D.~Brouer, D.~Borkmann, J.~Fastabend, T.~Herbert, D.~Ahern, D.~Miller.
\emph{The eXpress Data Path: Fast Programmable Packet Processing in the Operating System Kernel}.
Proceedings of the 14th International Conference on emerging Networking EXperiments and Technologies (CoNEXT), 2018.

\bibitem{heiser2018}
G.~Heiser et al.
\emph{Benchmarking Crimes: An Emerging Threat in Systems Security}.
arXiv:1801.02381, 2018.

\bibitem{lisa21}
Z.~Jones.
\emph{Performance Analysis of XDP Programs}.
USENIX LISA21, 2021.

\bibitem{cnnkernel2025}
Y.~Zheng, J.~Zhang.
\emph{In-Kernel CNN Inference for Edge Devices: An eBPF-Based Approach to Low-Latency and Resource-Efficient Processing}.
ICA3PP 2025, Lecture Notes in Computer Science, vol.~16387, Springer.
DOI: 10.1007/978-981-95-8417-8\_34.

\bibitem{nntree2024}
T.~Hara, M.~Sasabe.
\emph{Practicality of in-kernel/user-space packet processing empowered by lightweight neural network and decision tree}.
Computer Networks, vol.~240, art.~110188, pp.~1--18, febbraio 2024.
DOI: 10.1016/j.comnet.2024.110188.

\bibitem{n3ic2022}
G.~Siracusano, S.~Galea, D.~Sanvito, M.~Malekzadeh, G.~Antichi, P.~Costa, H.~Haddadi, R.~Bifulco.
\emph{Re-architecting Traffic Analysis with Neural Network Interface Cards}.
Proceedings of the 19th USENIX Symposium on Networked Systems Design and Implementation (NSDI), pp.~513--525, 2022.

\bibitem{inmlsurvey2024}
C.~Zheng, X.~Hong, D.~Ding, S.~Vargaftik, Y.~Ben-Itzhak, N.~Zilberman.
\emph{In-Network Machine Learning Using Programmable Network Devices: A Survey}.
IEEE Communications Surveys \& Tutorials, vol.~26, n.~2, pp.~1171--1200, 2024.
DOI: 10.1109/COMST.2023.3344351.

\end{thebibliography}

\end{document}
